\author{Kateryna Iefremenko} %use this field for editors as well
\title{The role of cross-linguistic influence, internal dynamics, and register in heritage speakers' languages}
\subtitle{Evidence from word order in heritage Turkish and Kurmanji Kurdish}
\ISBNdigital{978-3-96110-582-3}
\ISBNhardcover{978-3-98554-199-7}
\BookDOI{10.5281/zenodo.20544154}
\typesetter{Sebastian Nordhoff}
\proofreader{Matthew Korte}
\lsCoverTitleSizes{51.5pt}{17pt}
\renewcommand{\lsSeries}{cib} 
\renewcommand{\lsID}{597}
\renewcommand{\lsSeriesNumber}{7}
\lsCoverTitleSizes{45pt}{15mm}


\renewcommand{\lsImpressumExtra}{Zugleich Promotion an der Universität Potsdam (2024). Gutachter Prof. Dr. Christoph Schroeder, Prof. Dr. Christine Czinglar \&  Prof. Dr. Geoffrey Haig. Datum der mündlichen Prüfung: 8.7.2024}


\BackBody{This monograph examines sources of variation in the grammars of heritage speakers, focusing on word order in Turkish and Kurmanji Kurdish as heritage languages in the contexts of majority German and majority Turkish. By bringing together the relatively well-studied case of heritage Turkish and the less extensively researched case of heritage Kurdish, the study develops a comparative framework for identifying mechanisms that shape heritage grammars. It evaluates several explanatory factors, including cross-linguistic influence, language-internal dynamics, and register. The analysis also moves beyond a monolingual–bilingual contrast by including bilingual speakers for whom Turkish serves as the societal majority language, thereby shifting the focus from bilingualism per se to the role of a language’s societal status as heritage or majority. In sum, the findings demonstrate that register is a central factor shaping variation in word order of heritage Turkish, whereas cross-linguistic influence plays a more prominent role in heritage Kurdish. More broadly, the study underscores the need to expand the empirical scope of heritage language research to include a wider range of languages and sociolinguistic settings, fostering closer dialogue with studies of minoritized languages and multilingual language communities.}
